\section{Geometric-phase foundations for the holonomy template}
\label{app:holonomy}

For a family of normalized states $\ket{\psi(\lambda)}$ parameterized by $\lambda$ on a manifold, define the Berry connection and curvature \cite{Simon1983,Berry1984,Nakahara2003}
\begin{equation}
A = \iu\,\braket{\psi|\diff\psi},\qquad F=\diff A.
\end{equation}
For a closed loop $C$, the geometric phase is
\begin{equation}
\gamma_{\mathrm{geom}}[C]=\oint_C A = \int_{S:\,\partial S=C} F,
\end{equation}
which is gauge-invariant at the level of curvature. The mixing-holonomy axiom in Axiom~\ref{ax:mixing_holonomy} uses the non-Abelian generalization with a $U(n)$ connection and path-ordered exponentiation \cite{WilczekZee1984,Nakahara2003}.

\begin{proposition}[Any unitary can be realized as holonomy on a circle]
\label{prop:any_unitary_holonomy}
For any $U\in U(n)$, there exists a $U(n)$ connection on the trivial bundle over $S^1$ whose holonomy along the generator of $\pi_1(S^1)$ equals $U$.
\end{proposition}
\begin{proof}
Choose a matrix logarithm $H$ such that $U=\exp(\iu H)$. On $S^1$ with angular coordinate $\theta\in[0,2\pi]$, define a constant $u(n)$-valued connection one-form $\mathcal{A}=\iu H\,\diff\theta/(2\pi)$. Then the path-ordered exponential reduces to an ordinary exponential and
\[
\mathcal{P}\exp\!\left(\oint_{S^1} \mathcal{A}\right)=\exp\!\left(\int_0^{2\pi} \frac{\iu H}{2\pi}\,\diff\theta\right)=\exp(\iu H)=U.
\]
Standard holonomy theory is developed in \cite{KobayashiNomizu1963}.
\end{proof}

\begin{remark}[Predictivity requires geometric rigidity]
\label{rem:holonomy_rigidity}
Proposition~\ref{prop:any_unitary_holonomy} shows that a holonomy representation is formally universal. Therefore, a predictive constant-geometry model must specify $\mathcal{M}$ and $\mathcal{A}$ by geometric rigidity (e.g., fixed bundles/metrics, minimal cycles, symmetry constraints), rather than by free fitting of a general connection.
\end{remark}
